\documentclass[11pt]{article}

\usepackage[T1]{fontenc}
\usepackage[utf8]{inputenc}
\usepackage{amsmath,amssymb,amsthm,mathtools}
\usepackage{microtype}
\usepackage{hyperref}
\usepackage{geometry}
\geometry{margin=1in}

\newtheorem{theorem}{Theorem}[section]
\newtheorem{proposition}[theorem]{Proposition}
\newtheorem{lemma}[theorem]{Lemma}
\newtheorem{corollary}[theorem]{Corollary}
\newtheorem{remark}[theorem]{Remark}
\newtheorem{definition}[theorem]{Definition}

\newcommand{\C}{\mathbb{C}}
\newcommand{\op}{\mathrm{op}}
\newcommand{\cmax}{C_{\max}}
\newcommand{\dd}{\delta}
\newcommand{\etaq}{\eta}
\newcommand{\tauq}{\tau}
\newcommand{\PhiRho}{\Phi_{\rho}}

\title{Sharp Quantitative Stability of Finite Projective Records\\
under Approximately Commuting Unitary Transformations}

\author{Bertrand Dalimier\\
Independent Researcher}

\date{}

\begin{document}

\maketitle

\begin{abstract}
Let $(P_i)_{i\in I}$ be a finite orthogonal projective decomposition of a
finite-dimensional complex Hilbert space and let $U$ be unitary.  Write
$a_i=\|[P_i,U]\|_{\op}$ for the cellwise commutator defects,
$\dd^2=\sum_i a_i^2$ for their global quadratic budget, and
\[
  \cmax=\max_{S\subseteq I}\|[P_S,U]\|_{\op},
  \qquad P_S=\sum_{i\in S}P_i ,
\]
for the finite cut envelope.  We prove the sharp universal estimate
$\cmax^2\le \dd^2/2$.  When $\dd>0$, equality holds if and only if exactly two
cells have nonzero commutator defect.  We then establish an inverse stability
theorem: if
\[
  \cmax^2\ge (1-\etaq)\frac{\dd^2}{2},\qquad
  0\le\etaq\le\frac15,\qquad \dd\ge\rho>0,
\]
then two cells capture all but a relative fraction
\[
  \PhiRho(\etaq)
  =2\etaq+\frac{36\etaq}{(1-3\etaq)\rho^2}
\]
of the global quadratic budget.  For $\etaq\le1/6$ this is bounded by
$\etaq(2+72/\rho^2)$.  Finally, an explicit three-cell family shows that the
inverse-square scale dependence is necessary up to constants: in an iterated
near-saturation regime, $(\tauq/\etaq)\dd^2\to4$.  All principal statements
and the sharpness construction have been formally verified in Lean 4.
\end{abstract}

\noindent\textbf{Keywords.}
projector commutators; quantitative stability; rigidity; near equality;
unitary operators; formal verification

\section{Introduction}

Finite projective decompositions provide a simple setting in which a unitary
transformation can be compared with a prescribed collection of mutually
orthogonal sectors.  If $P_i$ is the projector associated with sector $i$,
then the commutator $[P_i,U]$ measures the failure of $U$ to preserve that
sector.  The individual operator norms
\[
  a_i:=\|[P_i,U]\|_{\op}
\]
therefore define a nonnegative finite defect profile.  The quadratic aggregate
\[
  \dd^2:=\sum_{i\in I} a_i^2
\]
is a natural dimension-free global measure of failure of simultaneous
commutation.

A second object is obtained by grouping sectors into a cut.  For
$S\subseteq I$, let
\[
  P_S:=\sum_{i\in S}P_i,\qquad
  C_S:=\|[P_S,U]\|_{\op}.
\]
Since $I$ is finite, the cut envelope
\[
  \cmax:=\max_{S\subseteq I}C_S
\]
is attained.  The basic inequality studied here is
\begin{equation}\label{eq:envelope}
  \cmax^2\le \frac{\dd^2}{2}.
\end{equation}
It is a finite operator-theoretic inequality: no probability assignment,
dynamical assumption, open-system structure, or asymptotic Hilbert-space
limit is involved.

The purpose of this paper is to analyze the near-equality regime of
\eqref{eq:envelope}.  There are three levels.

First, we identify the exact equality cases.  For positive global defect,
saturation of \eqref{eq:envelope} is equivalent to the statement that exactly
two cells are active: all other cellwise commutator defects vanish.

Second, we prove a quantitative inverse theorem.  If the cut envelope nearly
saturates \eqref{eq:envelope}, then almost all of the quadratic defect budget
must be carried by two cells.  The estimate is independent of the number of
cells.  At a fixed positive defect scale $\rho$, the relative two-cell tail is
$O(\etaq/\rho^2)$ as the saturation deficit $\etaq$ tends to zero.

Third, we construct an explicit three-cell family showing that the
$1/\rho^2$ scale is not a proof artifact.  The family has closed formulas for
the three cell budgets, the global defect, the maximal cut, the
near-saturation parameter, and the two-cell tail.  Its asymptotics give the
constant
\[
  \lim_{m\to\infty}\left[
    \lim_{n\to\infty}\frac{\tauq_{n,m}}{\etaq_{n,m}}\dd_{n,m}^2
  \right]=4.
\]
Thus a stability estimate uniform down to zero defect cannot in general have
a tail bound of order $O(\etaq)$ with a scale-independent constant.

The arguments separate into a finite real concentration layer and an
operator-theoretic layer.  This separation is also reflected in a companion
Lean 4 formalization.  The principal theorem statements, equality
classification, rank-one formula, explicit sharpness family, and asymptotic
limit are machine checked.  The formal development is used here as a
verification artifact rather than as the mathematical subject of the paper.

% TODO before submission:
% 1. Complete systematic priority/related-work audit.
% 2. Replace repository commit by archival release/DOI.
% 3. Adapt notation/template to target journal.

\section{Finite projective setup}

Let $H$ be a finite-dimensional complex Hilbert space.  Let
$I$ be a finite index set and let $(P_i)_{i\in I}$ be orthogonal projections
satisfying
\begin{equation}\label{eq:decomp}
  P_iP_j=0\quad(i\neq j),
  \qquad
  \sum_{i\in I}P_i=I_H.
\end{equation}
We call the family a finite projective record or finite projective
decomposition.  Let $U:H\to H$ be unitary.

For each cell $i$, define
\begin{equation}\label{eq:ai}
  a_i:=\|[P_i,U]\|_{\op},
  \qquad
  p_i:=a_i^2,
\end{equation}
and define the global quadratic commutator defect
\begin{equation}\label{eq:delta}
  \dd:=\left(\sum_{i\in I}p_i\right)^{1/2}.
\end{equation}
For a cut $S\subseteq I$, set
\begin{equation}\label{eq:cut}
  P_S:=\sum_{i\in S}P_i,
  \qquad
  C_S:=\|[P_S,U]\|_{\op},
  \qquad
  \cmax:=\max_{S\subseteq I}C_S.
\end{equation}

The quantity $\dd$ is the $\ell^2$ norm of the cellwise operator-norm profile.
The quantity $\cmax$ instead probes the largest coherent defect visible after
aggregating an arbitrary subset of the cells.

For later use, when $\dd>0$ define the optimal two-cell tail
\begin{equation}\label{eq:tau}
  \tauq:=
  \min_{\substack{i,j\in I\\i\neq j}}
  \frac{\dd^2-p_i-p_j}{\dd^2}.
\end{equation}
Equivalently, $\tauq$ is the fraction of the total quadratic budget left after
retaining the best pair of distinct cells.

\begin{remark}
The formal theorem is stated without division by $\dd^2$: it asserts the
existence of distinct $i,j$ such that
\[
  \dd^2-p_i-p_j\le \varepsilon\,\dd^2.
\]
This formulation remains meaningful at zero defect.  In the positive-defect
regime it is equivalent to $\tauq\le\varepsilon$.
\end{remark}

\section{The finite cut envelope}

The key observation is that a cut commutator is controlled simultaneously by
the defect budget on the cut and by the defect budget on its complement.

Relative to the orthogonal splitting
\[
  H=P_SH\oplus (I-P_S)H,
\]
the commutator $[P_S,U]$ is off diagonal.  Its two cross blocks are
\[
  P_SU(I-P_S),
  \qquad
  (I-P_S)UP_S.
\]
For a unitary $U$ the operator norm of the cut commutator is the maximum of
the norms of these two cross blocks.  The incoming block is controlled by the
cellwise budgets on $S$, while the outgoing block is controlled by the
budgets on $S^c$.  Consequently,
\begin{equation}\label{eq:cut-min}
  C_S^2
  \le
  \min\left\{
    \sum_{i\in S}p_i,\,
    \sum_{i\notin S}p_i
  \right\}.
\end{equation}

Since the two quantities inside the minimum sum to $\dd^2$, we obtain the
universal envelope.

\begin{theorem}[Cut-envelope inequality]\label{thm:envelope}
For every finite projective decomposition and every unitary $U$,
\[
  \boxed{\cmax^2\le\frac{\dd^2}{2}}.
\]
\end{theorem}

\begin{proof}
For every $S\subseteq I$, \eqref{eq:cut-min} gives
\[
  C_S^2
  \le
  \min\{A_S,\dd^2-A_S\},
  \qquad
  A_S:=\sum_{i\in S}p_i.
\]
The elementary inequality
$\min\{x,T-x\}\le T/2$ for $0\le x\le T$ yields
$C_S^2\le\dd^2/2$.  Taking the maximum over the finite set of cuts proves the
claim.
\end{proof}

The coefficient $1/2$ is optimal.  More importantly for the present paper,
its equality cases are rigid.

\section{Exact rigidity at saturation}

Call a cell \emph{active} when $p_i>0$.

\begin{theorem}[Exact equality classification]\label{thm:rigidity}
Assume $\dd>0$.  Then
\[
  \boxed{
    \cmax^2=\frac{\dd^2}{2}
    \quad\Longleftrightarrow\quad
    \text{exactly two cells are active}.
  }
\]
More explicitly, equality holds if and only if there exist distinct
$i,j\in I$ such that
\[
  p_i>0,\qquad p_j>0,\qquad
  p_k=0\quad\text{for every }k\notin\{i,j\}.
\]
\end{theorem}

\paragraph{Proof architecture.}
Choose a cut $S$ attaining $\cmax$.  Equality in
Theorem~\ref{thm:envelope} forces equality throughout the two-sided budget
bound:
\[
  C_S^2
  =
  \sum_{i\in S}p_i
  =
  \sum_{i\notin S}p_i
  =
  \frac{\dd^2}{2}.
\]
The equality analysis of the cross-block estimates implies that the active
budget on each side of the cut is concentrated on a single cell.  Hence the
global support is contained in two cells, one in $S$ and one in $S^c$.
Because $\dd>0$ and both sides carry $\dd^2/2$, both cells are active.

Conversely, if the commutator profile is supported on two active cells,
choosing a cut that separates them makes both complementary budgets equal to
$\dd^2/2$ and saturates the cut estimate.

The positive-defect hypothesis is only used to exclude the degenerate case
in which every cell commutes with $U$.

\section{Quantitative stability of near-saturators}

We now quantify Theorem~\ref{thm:rigidity}.  Let
$0\le\etaq\le1/5$ and assume
\begin{equation}\label{eq:near}
  \cmax^2
  \ge
  (1-\etaq)\frac{\dd^2}{2}.
\end{equation}
Equivalently,
\[
  2(1-\etaq)\dd^2\le (2\cmax)^2.
\]

A scale condition is necessary.  We therefore assume
\begin{equation}\label{eq:scale}
  \dd\ge\rho>0.
\end{equation}
Define
\begin{equation}\label{eq:modulus}
  \PhiRho(\etaq):=
  2\etaq+
  \frac{36\etaq}{(1-3\etaq)\rho^2}.
\end{equation}

\begin{theorem}[Quantitative two-cell stability]\label{thm:stability}
Suppose \eqref{eq:near} and \eqref{eq:scale} hold with
$0\le\etaq\le1/5$.  Then there exist distinct cells $i,j$ such that
\begin{equation}\label{eq:stability-conclusion}
  \boxed{
  \dd^2-p_i-p_j
  \le
  \PhiRho(\etaq)\,\dd^2.
  }
\end{equation}
Equivalently,
\[
  \tauq\le
  2\etaq+
  \frac{36\etaq}{(1-3\etaq)\rho^2}.
\]
\end{theorem}

At a fixed positive scale, this yields a genuine stability modulus:
\[
  \PhiRho(\etaq)\longrightarrow0
  \qquad(\etaq\downarrow0).
\]

\begin{corollary}[Transparent linear form]\label{cor:linear}
If in addition $\etaq\le1/6$, then
\begin{equation}\label{eq:linear}
  \boxed{
    \tauq\le
    \etaq\left(2+\frac{72}{\rho^2}\right).
  }
\end{equation}
\end{corollary}

\subsection{Proof mechanism}

The proof is most transparent when one fixes a maximizing cut $S$ and writes
\[
  A:=\dd^2,\qquad
  C:=C_S=\cmax,\qquad
  E:=A-2C^2.
\]
Near saturation gives
\[
  0\le E\le\etaq A.
\]
The cross-block representation supplies a unit norm witness realizing the
operator norm of one side of the cut.  Its cellwise decomposition produces a
finite nonnegative profile.  A normalized Gram estimate bounds the
off-diagonal spread of that profile in terms of $E$.  A purely finite
concentration lemma then gives a dominant cell on that side of the cut.

The same argument is applied to the inverse unitary, which exchanges incoming
and outgoing blocks and therefore gives a dominant cell on the complementary
side.  The two witnesses are necessarily distinct.  Combining the bad-cell
mass and the two one-sided tails gives an estimate of the form
\[
  \dd^2-p_i-p_j
  \le
  2E+\frac{18E}{d},
\]
where the denominator $d$ is bounded below by
\[
  d\ge\frac{(1-3\etaq)A}{2}.
\]
Using $E\le\etaq A$ and $A=\dd^2\ge\rho^2$ yields
\[
  \frac{\dd^2-p_i-p_j}{\dd^2}
  \le
  2\etaq+
  \frac{36\etaq}{(1-3\etaq)\rho^2}.
\]
The constants here are explicit rather than optimized.  Section~\ref{sec:sharp}
shows, however, that the inverse-square dependence on the scale cannot in
general be removed.

\section{An explicit three-cell sharpness family}\label{sec:sharp}

We construct a family in dimension three whose near-saturation deficit tends
to zero while the ratio $\tauq/\etaq$ diverges at the inverse-square defect
scale.

Let $(p_n,q_n)$ be the positive integer recurrence
\[
  (p_0,q_0)=(1,1),
\]
\[
  p_{n+1}=3p_n+4q_n,
  \qquad
  q_{n+1}=2p_n+3q_n.
\]
It preserves
\begin{equation}\label{eq:pell}
  p_n^2+1=2q_n^2.
\end{equation}
Define
\[
  a_n=\frac{p_n+1}{2q_n},
  \qquad
  b_n=\frac{p_n-1}{2q_n},
\]
so that $a_n^2+b_n^2=1$.  Rotate these amplitudes by $45^\circ$:
\begin{equation}\label{eq:AB}
  A_n=\frac{a_n+b_n}{\sqrt2},
  \qquad
  B_n=\frac{a_n-b_n}{\sqrt2}.
\end{equation}
Then
\[
  A_n^2+B_n^2=1,\qquad
  B_n=\frac{1}{\sqrt2\,q_n},
\]
hence
\begin{equation}\label{eq:AB-limits}
  B_n^2\to0,\qquad A_n^2\to1.
\end{equation}

For a second index $m$, let $c_m,s_m$ be the explicit rational unit-circle
parameters
\[
  c_m^2+s_m^2=1,\qquad s_m>0,
\]
coming from the same Pell family, and set
\begin{equation}\label{eq:d}
  d_m:=1-c_m.
\end{equation}
Thus
\begin{equation}\label{eq:s-d}
  s_m^2=2d_m-d_m^2,
  \qquad
  d_m\to0.
\end{equation}

The three-dimensional unitary is obtained by conjugating the planar rotation
with a basis change whose first direction is
$(A_n,B_n,0)$.  For the three rank-one cells, the exact squared commutator
budgets are
\begin{align}
  p_0 &=
  2d_mA_n^2-d_m^2A_n^4, \label{eq:p0}\\
  p_1 &=
  2d_mB_n^2-d_m^2B_n^4, \label{eq:p1}\\
  p_2 &= s_m^2. \label{eq:p2}
\end{align}
These identities follow from the rank-one formula
\begin{equation}\label{eq:rankone}
  \boxed{
  \|[P_e,U]\|_{\op}^2
  =
  1-\bigl|\langle e,Ue\rangle\bigr|^2
  }
\end{equation}
for a unit vector $e$ and the rank-one projector $P_e$.

Summing \eqref{eq:p0}--\eqref{eq:p2} gives
\begin{equation}\label{eq:global-sharp}
  \boxed{
  \dd_{n,m}^2
  =
  2s_m^2+
  2d_m^2A_n^2B_n^2.
  }
\end{equation}
The maximal cut is exact:
\begin{equation}\label{eq:cmax-sharp}
  \boxed{
  \cmax(n,m)=s_m.
  }
\end{equation}

Define the relative saturation deficit
\[
  \etaq_{n,m}
  :=
  \frac{\dd_{n,m}^2-2\cmax(n,m)^2}{\dd_{n,m}^2}.
\]
Then
\begin{equation}\label{eq:eta-sharp}
  \boxed{
  \etaq_{n,m}
  =
  \frac{
    2d_m^2A_n^2B_n^2
  }{
    2s_m^2+2d_m^2A_n^2B_n^2
  }.
  }
\end{equation}
Retaining cells $0$ and $2$ leaves the cell-$1$ tail
\begin{equation}\label{eq:tau-sharp}
  \tauq_{n,m}
  =
  \frac{
    2d_mB_n^2-d_m^2B_n^4
  }{
    2s_m^2+2d_m^2A_n^2B_n^2
  }.
\end{equation}
The exact quotient simplifies to
\begin{equation}\label{eq:ratio-sharp}
  \boxed{
  \frac{\tauq_{n,m}}{\etaq_{n,m}}
  =
  \frac{2-d_mB_n^2}{2d_mA_n^2}.
  }
\end{equation}

\begin{theorem}[Scale sharpness]\label{thm:scale-sharp}
For fixed $m$,
\begin{equation}\label{eq:fixed-m-ratio}
  \lim_{n\to\infty}
  \frac{\tauq_{n,m}}{\etaq_{n,m}}
  =
  \frac1{d_m},
\end{equation}
and
\begin{equation}\label{eq:fixed-m-delta}
  \lim_{n\to\infty}\dd_{n,m}^2=2s_m^2.
\end{equation}
Consequently,
\begin{equation}\label{eq:iterated4}
  \boxed{
  \lim_{m\to\infty}
  \left[
    \lim_{n\to\infty}
    \frac{\tauq_{n,m}}{\etaq_{n,m}}
    \dd_{n,m}^2
  \right]
  =4.
  }
\end{equation}
\end{theorem}

\begin{proof}
Equations \eqref{eq:fixed-m-ratio} and \eqref{eq:fixed-m-delta} follow from
\eqref{eq:AB-limits}, \eqref{eq:global-sharp}, and
\eqref{eq:ratio-sharp}.  Their product tends, for fixed $m$, to
\[
  \frac{2s_m^2}{d_m}.
\]
Using \eqref{eq:s-d},
\[
  \frac{2s_m^2}{d_m}
  =
  4-2d_m.
\]
Finally $d_m\to0$, giving \eqref{eq:iterated4}.
\end{proof}

\begin{corollary}[Necessity of the inverse-square scale]
Along the family above,
\[
  \frac{\tauq_{n,m}}{\etaq_{n,m}}
  \asymp \frac{1}{\dd_{n,m}^2}
\]
in the iterated small-defect regime.  Hence a dimension-free near-equality
theorem with tail $O(\etaq)$ and a constant independent of the defect scale
cannot hold in general.  The power $1/\rho^2$ in
Theorem~\ref{thm:stability} is therefore necessary up to multiplicative
constants.
\end{corollary}

\section{Relation with the sharp record-transfer bound}

The cut envelope also enters the sharp uniform transfer inequality for finite
projective record profiles.  For a normalized state $x$, the companion
formal development proves
\begin{equation}\label{eq:transfer}
  \|r_D(Ux)-r_D(x)\|_1
  \le
  \sqrt2\,\dd,
\end{equation}
with the universal coefficient $\sqrt2$ optimal.

The present paper addresses the inverse geometry behind near saturation of
the cut stage that feeds \eqref{eq:transfer}.  Theorems
\ref{thm:rigidity} and \ref{thm:stability} show that extremal and
near-extremal cut behavior is effectively two-cell.  This part of the
argument is independent of the Born-rule interpretation of the record
profile: it depends only on finite projectors, unitary maps, operator norms,
and finite real concentration.

A possible extension is to combine the present cut-envelope stability with a
separate stability analysis of every remaining inequality in the derivation
of \eqref{eq:transfer}.  Such a result would classify near-extremizers of the
full $\sqrt2$ transfer inequality, rather than only the operator-theoretic cut
stage.  We do not claim such a classification here.

\section{Formal verification}

The mathematical core has a companion Lean 4 formalization in the repository
\[
  \texttt{Bobart0/everettian-decoherence-lean}.
\]
The development is organized so that the publication-facing surface imports
the exact rigidity theorem, the quantitative stability theorem, and the
sharpness asymptotics.  A dedicated audit module checks the principal
declarations and prints their axiomatic dependencies.

At manuscript version 0.1, the reference development is the CI-verified
repository state containing the following principal declarations:
\begin{center}
\begin{tabular}{p{0.39\textwidth}p{0.52\textwidth}}
\hline
Mathematical role & Lean declaration\\
\hline
Cut envelope &
\texttt{maxSubsetCommutatorOpNorm\_sq\_le\_half\_globalSq}\\
Exact rigidity &
\texttt{maxCut\_saturation\_iff\_exactlyTwoActiveCellCommutators}\\
Quantitative modulus &
\texttt{twoCellCommutatorConcentratedWithin\_of\_maxCut\_near\_saturated}\\
Linear corollary &
\texttt{twoCellCommutatorConcentratedWithin\_linear\_of\_maxCut\_near\_saturated}\\
Rank-one identity &
\texttt{perspectiveProjectorCommutatorOpNormProfile\_sq\_rankOne}\\
Sharpness ratio &
\texttt{stabilitySharpnessTau\_div\_Eta\_exact}\\
Fixed-scale asymptotics &
\texttt{stabilitySharpnessTau\_div\_Eta\_tendsto\_inv\_D}\\
Defect-scale asymptotics &
\texttt{stabilitySharpness\_delta\_sq\_tendsto\_two\_s\_sq}\\
Constant-$4$ obstruction &
\texttt{stabilitySharpnessScaleLimit\_tendsto\_four}\\
\hline
\end{tabular}
\end{center}

The publication-facing facade is
\texttt{EverettianDecoherence.Approximation.StabilitySharpness}; the dedicated
audit module is
\texttt{EverettianDecoherence.Audit.StabilityRigiditySharpness}.
The final formalization head preceding this manuscript branch is
\texttt{81eb3ff53f6219749b84491513ff437fe7c545c1}, whose complete CI passed,
including the library build, audit aggregate, guards, and diff-hygiene checks.

For archival submission, this commit reference should be replaced or
supplemented by a tagged release and a persistent archive identifier.

\section{Discussion and limitations}

The stability theorem is dimension-free in the number of projective cells:
its constants do not grow with $|I|$.  The price is the explicit positive
scale $\rho$.  The sharpness family shows that this dependence is structural
rather than merely technical.

The theorem is also deliberately algebraic.  It does not assert that the
projective decomposition is dynamically selected, that the sectors are
environmentally redundant, that off-diagonal density-matrix elements decay
in time, or that a preferred basis emerges.  Nor does the operator-theoretic
core require a probability interpretation of the cells.  Such physical
questions require additional hypotheses that are outside the scope of the
present result.

The constants $36$ and $72$ arise from the quantitative proof architecture
and are not claimed to be optimal.  The sharpness construction identifies the
correct inverse-square scale and an asymptotic constant $4$ for the explicit
family, but it does not by itself determine the best universal coefficient in
the full finite-dimensional stability theorem.

\section{Related work and priority status}

The operator-theoretic background includes several neighboring but distinct
stability questions.  Halmos' analysis of pairs of subspaces provides a
canonical geometric framework for two projections
\cite{Halmos1969TwoSubspaces}, while Davis--Kahan perturbation theory
quantifies the displacement of invariant subspaces under spectral
perturbations \cite{DavisKahan1970}.  These results concern the geometry of
subspaces or spectral projectors, rather than the finite cut envelope studied
here.

A second neighboring line asks whether matrices with small commutators are
close to exactly commuting matrices.  For self-adjoint matrices, Lin proved
dimension-independent approximation by commuting self-adjoints
\cite{Lin1997AlmostCommuting}; see also the short proof of Friis and R{\o}rdam
\cite{FriisRordam1996} and quantitative work of Hastings
\cite{Hastings2009}.  The unitary problem has a different obstruction theory:
Exel and Loring exhibited and analyzed topological obstructions for almost
commuting unitaries \cite{ExelLoring1989,ExelLoring1991}, and effective
commutator-control questions remain active
\cite{DorOnHallKachkovskiy2025}.

The present problem is not an approximation-to-a-commuting-tuple theorem.
The projectors $(P_i)$ are fixed and already commute mutually; the unitary
$U$ is not replaced by a nearby commuting operator.  Instead, we study a
specific finite functional of the family of commutators $[P_i,U]$: the
largest operator norm after aggregating cells into an arbitrary cut.  The
main questions are the equality structure of its sharp quadratic envelope
and the stability of its near-equality cases.

The searches completed for this draft have not identified a prior theorem
with the same combination of (i) the finite cut envelope
$\cmax^2\le\dd^2/2$, (ii) an exact two-active-cell equality classification,
(iii) a dimension-free near-saturation theorem with an explicit scale
parameter, and (iv) a three-cell family proving the necessity of the
inverse-square scale.  This is not yet a claim of historical priority.
Before submission, the search will be extended systematically through
operator-theory databases and citation chains, with particular attention to
near-equality results for projection commutators and block-operator norm
inequalities.

\section{Conclusion}

For a finite projective decomposition and a unitary transformation, the
largest aggregate cut commutator satisfies the sharp quadratic envelope
\[
  \cmax^2\le\frac{\dd^2}{2}.
\]
Positive-defect equality is rigid: it occurs exactly when two cells, and only
two cells, carry nonzero commutator defect.  Near equality inherits a
dimension-free quantitative form.  At defect scale at least $\rho$,
near-saturation with deficit $\etaq$ forces all but
\[
  2\etaq+
  \frac{36\etaq}{(1-3\etaq)\rho^2}
\]
of the relative quadratic budget onto two cells.

An explicit three-cell family proves that the inverse-square scale cannot be
eliminated: the asymptotic scaled ratio tends to $4$.  This closes the
qualitative--quantitative--sharpness chain for the finite cut envelope and
provides a machine-checked basis for further study of near-extremizers of
projective record-transfer inequalities.

\bibliographystyle{plain}
\bibliography{references_v01}

\end{document}
